Magnetic microscopy is rapidly advancing toward higher spatial resolutions and smaller sensor-to-sample distances. Under these conditions, dipolar approaches become increasingly limited because higher-order magnetic contributions can no longer be neglected when retrieving reliable full-vector information for paleomagnetic and rock magnetic applications. However, multipolar representations often require prior knowledge of source geometry and expansion center location, which can involve costly auxiliary measurements and additional processing complexity. This challenge becomes particularly relevant for planar magnetic microscopy observations, where the incomplete sampling geometry violates the ideal assumptions of spherical harmonic (SH) expansions. Here, we present an Euler-based approach that provides initial estimates of the expansion center for SH inversions, with optional nonlinear refinement through simplex optimization. The proposed methodology was evaluated using both synthetic micromagnetic models and real magnetic microscopy data, and its performance was compared with conventional dipolar approaches. Synthetic experiments with known ground-truth models demonstrated improved accuracy of SH-based inversions, particularly for near-surface measurements where higher-order magnetic signals become significant, while maintaining comparable performance at larger sensor heights. Application to real data acquired at approximately $5~\mu m$ sensor height showed broadly similar results for both approaches in most cases, although the SH framework produced more coherent solutions for anomalies affected by external magnetic contributions. These results represent an important step toward retrieving complex magnetic source information directly from planar magnetic anomaly measurements.